%==============================================================
% ΚΕΦΑΛΑΙΟ 4 — ΠΕΙΡΑΜΑΤΙΚΑ ΑΠΟΤΕΛΕΣΜΑΤΑ
%==============================================================
\chapter{Πειραματικά Αποτελέσματα}
\label{ch:results}

Στο παρόν κεφάλαιο παρουσιάζονται τα αποτελέσματα που προκύπτουν
από τα δεδομένα, τα μοντέλα και τις μεθοδολογίες που περιεγράφηκαν
αναλυτικά στο Κεφάλαιο~\ref{ch:methodology}.
Τα αποτελέσματα εστιάζουν στους τέσσερις άξονες: απόδοση,
ερμηνευσιμότητα, μεροληψία και παρακολούθηση.
Η σειρά παρουσίασης τους είναι ίδια με αυτή της μεθοδολογίας.
Το κύριο μοντέλο της εργασίας είναι το XGBoost, ενώ παράλληλα
δοκιμάζονται και τα μοντέλα Logistic Regression ως βασικό εργαλείο
του πιστωτικού κίνδυνου, αλλά και το FFNN ως πολυπλοκότερο μοντέλο
τεχνητής νοημοσύνης.
Για κάθε αποτέλεσμα υπάρχει και ένα τουλάχιστον γράφημα που βοηθάει
στην οπτικοποίηση του, ενώ παράλληλα παρέχεται ανάλυση ως προς
την σημαντικότητα και ερμηνεία του.
Κεντρικό εύρημα του κεφαλαίου αποτελεί η μη γραμμική σχέση ανάμεσα
στην πολυπλοκότητα, την απόδοση και τις απαιτήσεις επικύρωσης.
Οι τρεις αυτοί άξονες έχουν αυξημένη βαρύτητα στην επιλογή μοντέλου
στον πιστωτικό κίνδυνο, γεγονός που δίνει πρακτική χρησιμότητα
στα ευρήματα.

%--------------------------------------------------------------
\section{Αξιολόγηση Απόδοσης Μοντέλων}
\label{sec:performance_results}

\subsection{Μετρικές Αξιολόγησης}
\label{subsec:metrics}
Η εκπαίδευση των μοντέλων έγινε με τις παραμέτρους που
περιγράφηκαν αναλυτικά στο Κεφάλαιο~\ref{ch:methodology}.
Η αξιολόγηση έγινε επάνω στο test set, μέρος της βάσης
HELOC που προέκυψε από την διαίρεση της.
Τα παρακάτω διαγράμματα των μετρικών AUC-ROC, Gini και
KS statistic αποκαλύπτουν την πλήρη εικόνα
όσον αφορά τον πρώτο άξονα αξιολόγησης τους, την απόδοση.
Για κάθε μετρική τα αποτελέσματα παρουσιάζονται με πλήρη
ανάλυση στη συνέχεια.



\subsection{Σύγκριση Μοντέλων και Κατανομή Σκορ}
\label{subsec:model_comparison}

Η παρούσα υποενότητα περιέχει την ανάλυση και τις τιμές
των μετρικών που αναφέρθηκαν στην
υποενότητα~\ref{subsec:metrics}.
Ο Πίνακας~\ref{tab:model_metrics} συνοψίζει τις τιμές
των AUC, Gini και KS για τα μοντέλα Logistic Regression,
XGBoost και FFNN.

\begin{table}[htbp]
    \centering
    \caption[Σύγκριση μετρικών απόδοσης — test set]{Σύγκριση μετρικών απόδοσης στο test set.
    Με έντονη γραφή η καλύτερη τιμή ανά μετρική.}
    \label{tab:model_metrics}
    \begin{tabular}{lccc}
        \toprule
        \textbf{Μοντέλο} & \textbf{AUC} & \textbf{Gini} & \textbf{KS} \\
        \midrule
        Logistic Regression & 0.789          & 0.579          & 0.455 \\
        XGBoost             & \textbf{0.796} & \textbf{0.593} & 0.458 \\
        FFNN                & 0.793          & 0.587          & \textbf{0.465} \\
        \bottomrule
    \end{tabular}
\end{table}


% -------------------------------------------------------------------
% 04_model_comparison.png
% Επεξήγηση ευρημάτων της 04_model_comparison.png
% -------------------------------------------------------------------

Όπως φαίνεται στο Σχήμα~\ref{fig:model_comparison},
στις μετρικές AUC-ROC (0{,}796) και Gini (0{,}593)
το μοντέλο XGBoost υπερτερεί έναντι των Logistic
Regression (0{,}789, 0{,}579) και του FFNN (0{,}793,
0{,}587).
Στη μετρική KS το μοντέλο FFNN (0{,}465) έχει ελαφρώς
καλύτερη απόδοση από το XGBoost (0{,}458) και
Logistic Regression (0{,}455).
Οι διαφορές μεταξύ των μοντέλων και στους τρεις
δείκτες είναι οριακές, οπότε δεν συνιστούν πραγματική
διαφορά στην απόδοση και αποτελεσματικότητα αυτών.
Οι μικρές αυτές διαφορές ήταν αναμενόμενες, καθώς
οφείλονται στη γραμμικότητα της βάσης FICO HELOC
όπως αναφέρθηκε στην ενότητα~\ref{sec:dataset}.
Το παρόν εύρημα υποστηρίζει την υπόθεση ότι η απόδοση
δεν είναι ανάλογη της πολυπλοκότητας του μοντέλου.

% -------------------------------------------------------------------
% 04_model_comparison.png
% Σχήμα 4.1 — Grouped bar chart σύγκρισης μοντέλων (NB02)
% Τοποθέτηση: αμέσως μετά τον Πίνακα 4.1
% -------------------------------------------------------------------
\begin{figure}[htbp]
    \centering
    \includegraphics[width=\textwidth]{04_model_comparison}
    \caption[Σύγκριση μοντέλων — AUC, Gini, KS]{Σύγκριση των τριών μοντέλων στο test set για τους δείκτες AUC-ROC,
             Gini και KS. Οι μικρές διαφορές αποδόσεως αναδεικνύουν ότι η
             αυξημένη πολυπλοκότητα δεν συνεπάγεται ανάλογη βελτίωση.}
    \label{fig:model_comparison}
\end{figure}

\subsection{Καμπύλες ROC}
\label{sec:roc}

Το Σχήμα~\ref{fig:roc_curves} παρουσιάζει τις καμπύλες ROC και των τριών μοντέλων
μαζί με τη διαγώνιο τυχαίας βασικής γραμμής (AUC = 0{,}500). Και τα τρία μοντέλα
βρίσκονται σαφώς πάνω από τη διαγώνιο, επιβεβαιώνοντας την προβλεπτική τους
ικανότητα. Η οπτική απόσταση μεταξύ τους είναι πολύ μικρή, συνεπής με τα
αριθμητικά αποτελέσματα του Πίνακα~\ref{tab:model_metrics}. Η καμπύλη ROC
είναι ανεξάρτητη από το επιλεγμένο threshold, άρα δεν επηρεάζεται από την
επιχειρηματική απόφαση για το όριο αποδοχής δανείου. Αυτό την καθιστά ιδανική
μετρική για τη σύγκριση της ικανότητας διαχωρισμού μεταξύ μοντέλων.

% -------------------------------------------------------------------
% Σχήμα 4.2 — ROC Curves (NB02)
% Τοποθέτηση: αμέσως μετά τον έλεγχο καμπυλών ROC
% -------------------------------------------------------------------
\begin{figure}[htbp]
    \centering
    \includegraphics[width=0.80\textwidth]{14_roc_curves}
    \caption[Καμπύλες ROC τριών μοντέλων]{Καμπύλες ROC και των τριών μοντέλων στο test set. Η διαγώνιος
             αντιστοιχεί σε τυχαίο μοντέλο χωρίς προβλεπτική ικανότητα
             (AUC = 0{,}500). Η μικρή οπτική απόσταση μεταξύ των καμπυλών
             αντικατοπτρίζει τις μικρές διαφορές απόδοσης.}
    \label{fig:roc_curves}
\end{figure}



% 16_score_distribution.png
Τα τρία διαφορετικά overlapping histograms στο
Σχήμα~\ref{fig:score_distribution} οπτικοποιούν την πιθανοτική
κατανομή των δανειοληπτών για κάθε μοντέλο.
Οι δανειολήπτες με δάνειο που κατέληξε σε αθέτηση
συμβολίζονται με κόκκινο ενώ οι υπόλοιποι με πράσινο.
Ένα μοντέλο με τέλεια προβλεπτική ικανότητα θα παρήγε
διάγραμμα με τις δύο κατανομές συγκεντρωμένες στο 0 (Good)
και 1 (Bad).
Τα αποτελέσματα της εργασίας διαφέρουν από μια τέτοια
κατανομή καθώς παρατηρούνται περιπτώσεις Good και Bad
σε όλο σχεδόν το φάσμα, γεγονός αναμενόμενο για πραγματικά
δεδομένα που εισάγουν τυχαιότητα.
Είναι εμφανής και στα τρία διαγράμματα η τάση των μοντέλων
να τοποθετούν τις τιμές Good κοντά στο 0 και Bad κοντά στο 1.
Αυτή η συμπεριφορά υποδηλώνει προβλεπτική ικανότητα και
συμφωνεί με τα αποτελέσματα των προηγούμενων σχημάτων.
Τέλος, υπάρχει επικάλυψη κατανομών η οποία υποδηλώνει
αβεβαιότητα των μοντέλων ως προς την ταξινόμηση η οποία
όμως είναι ελεγχόμενη και αναμενόμενη σε μη ιδανικά μοντέλα.
Η παρατηρούμενη επικάλυψη είναι συνεπής με τις οριακές
διαφορές AUC και Gini του Πίνακα~\ref{tab:model_metrics},
επιβεβαιώνοντας ότι και οι τρεις μέθοδοι αξιολόγησης ---
πίνακας μετρικών, καμπύλες ROC και κατανομή σκορ ---
οδηγούν στο ίδιο συμπέρασμα.

% -------------------------------------------------------------------
% Σχήμα 4.3 — Κατανομή σκορ Bad/Good ανά μοντέλο (NB03)
% Τοποθέτηση: αμέσως μετά το κείμενο ανάλυσης
% -------------------------------------------------------------------
\begin{figure}[htbp]
    \centering
    \includegraphics[width=\textwidth]{16_score_distribution}
    \caption[Κατανομή προβλεπόμενων πιθανοτήτων ανά μοντέλο]{Κατανομή προβλεπόμενων πιθανοτήτων για Bad
    (κόκκινο) και Good (πράσινο) δανειολήπτες για τα τρία
    μοντέλα. Η επικάλυψη των δύο κατανομών αντικατοπτρίζει
    την αβεβαιότητα ταξινόμησης και είναι συνεπής με τις
    οριακές διαφορές AUC των μοντέλων.}
    \label{fig:score_distribution}
\end{figure}


%--------------------------------------------------------------
\subsection{Έλεγχος Υπερπροσαρμογής}
\label{subsec:overfitting}

Ο Πίνακας~\ref{tab:overfitting} περιέχει την απόδοση των
μοντέλων στο training set, στο validation set αλλά και το
ποσοστό της σχετικής του πτώσης.
Το ποσοστό αυτό αποτελεί βασικό εργαλείο για τον έλεγχο
υπερπροσαρμογής (overfitting), ενώ το όριο που
χρησιμοποιείται στην παρούσα εργασία είναι το 5\%.
Η βιβλιογραφία δεν ορίζει αυστηρό κατώφλι· τιμές μεταξύ 5\%--10\% θεωρούνται 
συνήθως πρακτικά αποδεκτές για μοντέλα με ικανότητα γενίκευσης 
\citep{cawley2010over}.
Το XGBoost είναι το μόνο μοντέλο που ξεπερνά ελαφρώς
το όριο.
Το αποτέλεσμα αυτό είναι αναμενόμενο καθώς εμφανίζεται
σε μεγάλο ποσοστό στα μοντέλα δένδρων αποφάσεων \citep{chen2016xgboost, hastie2009}.
Οι τεχνικές αντιμετώπισης που εφαρμόστηκαν ---
Early stopping και stochastic regularization μέσω
δειγματοληψίας δεδομένων και χαρακτηριστικών ---
είναι αρκετές για την αντιμετώπισή του καθώς η απόλυτη
διαφορά στην απόδοση ανάμεσα στα δύο sets είναι
μικρή (0{,}806 vs 0{,}796).
Τα υπόλοιπα μοντέλα βρίσκονται αρκετά κάτω του ορίου
με σχετικές τιμές 1{,}3\% και 2{,}6\%.
Συμπερασματικά, τα μοντέλα έχουν την ικανότητα να
γενικοποιούν τις πληροφορίες του training set και να
τις εφαρμόζουν ικανοποιητικά σε διαφορετικά δεδομένα.

\begin{table}[htbp]
    \centering
    \caption[Έλεγχος υπερπροσαρμογής — Train/Val AUC]{Έλεγχος υπερπροσαρμογής --- σύγκριση Train/Val AUC
    και σχετική πτώση ανά μοντέλο.}
    \label{tab:overfitting}
    \begin{tabular}{lcccc}
        \toprule
        \textbf{Μοντέλο} & \textbf{Train AUC} & \textbf{Val AUC}
        & \textbf{Διαφορά} & \textbf{Σχετική Πτώση} \\
        \midrule
        Logistic Regression & 0{,}802 & 0{,}791 & 0{,}011 & 1{,}3\% \\
        XGBoost             & 0{,}859 & 0{,}806 & 0{,}053 & 6{,}2\% \\
        FFNN                & 0{,}819 & 0{,}798 & 0{,}021 & 2{,}6\% \\
        \bottomrule
    \end{tabular}
\end{table}

%--------------------------------------------------------------
\subsection{Πίνακες Σύγχυσης και Ανάλυση Ορίου}
\label{subsec:confusion_threshold}

% 15_confusion_matrices.png
Οι Πίνακες Σύγχυσης του Σχήματος~\ref{fig:confusion_matrices}
παρουσιάζουν τις κατηγοριοποιήσεις TP, TN, FP και FN που
περιεγράφηκαν στην ενότητα~\ref{subsec:metrics} για κάθε
μοντέλο για προεπιλεγμένο threshold $0{,}5$.
Στον πιστωτικό κίνδυνο υπάρχει ασυμμετρία κόστους ανάμεσα
στα λάθη FN και FP με τα πρώτα να κοστίζουν περισσότερο
στους δανειοδότες.
Στη συγκεκριμένη τιμή threshold το μοντέλο Logistic Regression
έχει τη χειρότερη απόδοση στην κατηγορία FN ενώ το XGBoost
την καλύτερη.
Το αποτέλεσμα αυτό δείχνει ελαφρώς καλύτερη απόδοση του
XGBoost αλλά δεν μπορεί να γενικοποιηθεί για όλες τις
τιμές threshold.

% -------------------------------------------------------------------
% Σχήμα 4.4 — Confusion Matrices τριών μοντέλων (NB02)
% Τοποθέτηση: αμέσως μετά το κείμενο ανάλυσης
% -------------------------------------------------------------------
\begin{figure}[htbp]
    \centering
    \includegraphics[width=\textwidth]{15_confusion_matrices}
    \caption[Πίνακες σύγχυσης τριών μοντέλων]{Πίνακες σύγχυσης για τα τρία μοντέλα στο test set
    με threshold $0{,}5$. Το XGBoost εμφανίζει τον μικρότερο
    αριθμό FN (190), γεγονός που αντιστοιχεί στο χαμηλότερο
    κόστος λάθους για τον δανειοδότη.}
    \label{fig:confusion_matrices}
\end{figure}

% 18_threshold_analysis.png
Η επιλογή του threshold αποτελεί επιχειρηματική απόφαση
καθώς πρέπει να είναι προσανατολισμένη με την πολιτική
της τράπεζας και το επιτρεπόμενο ρίσκο που διατίθεται
να αναλάβει.
Στο Σχήμα~\ref{fig:threshold_analysis} φαίνονται οι τιμές
Precision (ποσοστό σωστά ταξινομημένων Bad από τους
προβλεπόμενους Bad) και Recall (ποσοστό Bad που εντοπίστηκαν
από το σύνολο των πραγματικών Bad) για όλες τις πιθανές
τιμές του threshold και για τα τρία μοντέλα.
Η μετρική F1 --- ο αρμονικός μέσος των Precision και
Recall --- δείχνει το βέλτιστο σημείο ισορροπίας.
Η διαφορά είναι αρκετή για να δώσει μια μικρή υπεροχή
στο μοντέλο XGBoost αλλά όχι να το καταστήσει σαφώς
ανώτερο.
Η ασυμμετρία κόστους πιθανώς να οδηγούσε μια τράπεζα
να διαλέξει μια τιμή κάτω του $0{,}5$ για μεγιστοποίηση
του Recall εις βάρος της Precision.

% -------------------------------------------------------------------
% Σχήμα 4.5 — Threshold Analysis τριών μοντέλων (NB02)
% Τοποθέτηση: αμέσως μετά το κείμενο ανάλυσης
% -------------------------------------------------------------------
\begin{figure}[htbp]
    \centering
    \includegraphics[width=\textwidth]{18_threshold_analysis}
    \caption[Threshold analysis τριών μοντέλων]{Επίδραση του threshold στις μετρικές Recall,
    Precision και F1 για τα τρία μοντέλα. Η κατακόρυφη
    διακεκομμένη γραμμή δείχνει το default threshold $0{,}5$.}
    \label{fig:threshold_analysis}
\end{figure}


%--------------------------------------------------------------
\subsection{Ανάλυση Κατά Δεκατημόρια}
\label{subsec:decile}

% 19_decile_analysis.png
Το Σχήμα~\ref{fig:decile_analysis} παρουσιάζει την ανάλυση
κατά δεκατημόρια για τα τρία μοντέλα.
Στο αριστερό μέρος του οπτικοποιείται το Bad Rate ανά
δεκατημόριο ενώ στο δεξί το αθροιστικό ποσοστό Bad
(Cumulative Bad Rate) σε σύγκριση με την τυχαία βασική
γραμμή, από την οποία τα μοντέλα σαφώς αποκλίνουν.
Οι τράπεζες απαιτούν στην decile analysis των αλγορίθμων
τα χειρότερα deciles να συγκεντρώνουν τον υψηλότερο κίνδυνο.
Η κατάταξη είναι ικανοποιητική με μη σημαντικές αριθμητικές διαφορές μεταξύ των αλγορίθμων.
Το εύρημα αυτό επιβεβαιώνει την προβλεπτική ικανότητα
των μοντέλων στον πιστωτικό κίνδυνο.

% -------------------------------------------------------------------
% Σχήμα 4.6 — Decile Analysis τριών μοντέλων (NB02)
% Τοποθέτηση: αμέσως μετά το κείμενο ανάλυσης
% -------------------------------------------------------------------
\begin{figure}[htbp]
    \centering
    \includegraphics[width=\textwidth]{19_decile_analysis}
    \caption[Decile analysis τριών μοντέλων]{Ανάλυση κατά δεκατημόρια για τα τρία μοντέλα.
    Αριστερά: Bad Rate ανά decile — τα πρώτα deciles
    συγκεντρώνουν τον υψηλότερο κίνδυνο.
    Δεξιά: Cumulative Bad Rate σε σύγκριση με τη διαγώνιο
    τυχαίας πρόβλεψης.}
    \label{fig:decile_analysis}
\end{figure}



\subsection{Βαθμονόμηση Μοντέλων}
\label{subsec:calibration}

% 22_calibration.png
Το Σχήμα~\ref{fig:calibration} παρουσιάζει τα διαγράμματα βαθμονόμησης
και των τριών μοντέλων.
Τα αποτελέσματα διαφέρουν από αυτά της μετρικής AUC καθώς εξετάζεται
κάθε αλγόριθμος μεμονωμένα σε κάθε διάστημα πιθανοτήτων.
Η διαγώνια διακεκομμένη γραμμή συμβολίζει την τέλεια βαθμονόμηση.
Όταν η απόδοση ενός μοντέλου την ξεπερνά αυτό μεταφράζεται σε
υπερεκτίμηση κινδύνου ενώ το αντίθετο σε υποεκτίμηση.
Συγκεκριμένα, η Logistic Regression ακολουθεί πιστά την διαγώνιο
σε όλο το εύρος τιμών με μια μικρή υπερεκτίμηση στα υψηλά επίπεδα
($>0{,}8$) η οποία είναι όμως αμελητέα.
XGBoost παρουσιάζει μια ζώνη υπερεκτίμησης κινδύνου στο διάστημα
$(0\text{--}0{,}4)$ και ένα διάστημα ηπιότερης υποεκτίμησης στο
διάστημα $(0{,}6\text{--}1)$.
Το διάγραμμα του FFNN εμφανίζει την πιο ακανόνιστη συμπεριφορά
από τα τρία με αιχμές, ενώ η εγκυρότητά του φαίνεται να είναι
καλύτερη στις ακραίες τιμές του συνόλου.
Τα αποτελέσματα συμφωνούν σε μεγάλο βαθμό με την βιβλιογραφία
στην οποία παρατηρείται τάση συμπίεσης πιθανοτήτων προς το κέντρο
για τα μοντέλα δένδρων~\citep{niculescu2005predicting}.
XGBoost εμφανίζει τάση υπερεκτίμησης κυρίως στα χαμηλά και
μεσαία επίπεδα.
Τέλος, η Logistic Regression είναι η πιο αξιόπιστα βαθμονομημένη,
γεγονός που εξηγείται από την εκπαίδευσή της με log-loss
(cross-entropy) που τιμωρεί τις αποκλίσεις πιθανότητας.
Έτσι καθίσταται καταλληλότερη για υπολογισμό PD στο πλαίσιο
Basel~\citep{basel2017}.

% -------------------------------------------------------------------
% Σχήμα 4.6 — Calibration Plots (NB02)
% Τοποθέτηση: αμέσως μετά την αναφορά στη βαθμονόμηση
% -------------------------------------------------------------------
\begin{figure}[htbp]
    \centering
    \includegraphics[width=\textwidth]{22_calibration}
    \caption[Διαγράμματα βαθμονόμησης τριών μοντέλων]{Διαγράμματα βαθμονόμησης των τριών μοντέλων. Η διαγώνιος
             αντιστοιχεί στην τέλεια αντιστοίχηση προβλεπόμενης πιθανότητας
             και παρατηρούμενου Bad Rate. Η βαθμονόμηση είναι απαραίτητη για
             την αξιόπιστη τιμολόγηση επιτοκίων και τον υπολογισμό εποπτικών
             κεφαλαίων (Basel III).}
    \label{fig:calibration}
\end{figure}



\section{Εξειδικευμένη Ανάλυση XGBoost}
\label{sec:results_xgboost}

\subsection{Ερμηνευσιμότητα SHAP --- Καθολικό Επίπεδο}
\label{sec:results_shap_global}

% 05_shap_feature_importance.png
Το Σχήμα~\ref{fig:shap_importance} παρουσιάζει το Global SHAP Feature Importance
για το XGBoost.
Όπως αναφέρθηκε στην ενότητα~\ref{sec:shap}, η Global SHAP ορίζεται ως η μέση
απόλυτη τιμή SHAP για κάθε χαρακτηριστικό σε ολόκληρο το test set.
Οι μεταβλητές με την μεγαλύτερη τιμή είναι αυτές που επηρεάζουν περισσότερο
τις προβλέψεις του μοντέλου συνολικά.
Συγκεκριμένα, κορυφαίο χαρακτηριστικό αναδεικνύεται το
\texttt{ExternalRiskEstimate} --- συνθετικό χαρακτηριστικό της FICO ---
το οποίο προκύπτει από την σύνθεση πολλαπλών πληροφοριών πιστωτικής
συμπεριφοράς σε έναν αριθμό.
Ακολουθεί το \texttt{MSinceMostRecentInqexcl7days} --- δείκτης πρόσφατων
πιστωτικών ερευνών --- το ποσοστό χρήσης ανακυκλούμενης πίστωσης
\texttt{NetFractionRevolvingBurden}, ο δείκτης ηλικίας πιστωτικού
ιστορικού \texttt{AverageMInFile} και το ποσοστό πιστωτικών λογαριασμών
χωρίς καμία καθυστέρηση πληρωμής \texttt{PercentTradesNeverDelinq}.
Τα αποτελέσματα αυτά αποκαλύπτουν ότι το μοντέλο θεωρεί πιο έμπιστους
τους δανειολήπτες που κυρίως δεν ψάχνουν για μεγάλο χρονικό διάστημα
ενεργά για δάνειο, δεν προσεγγίζουν το όριο της πιστωτικής τους κάρτας,
έχουν παλαιό πιστωτικό ιστορικό και χωρίς καθυστέρηση στις πληρωμές τους.
Λιγότερη βαρύτητα έχουν οι μεταβλητές \texttt{MSinceMostRecentTradeOpen}
και \texttt{NumTotalTrades}, διότι χαρακτηριστικά όπως οι μήνες από το
άνοιγμα του πιο πρόσφατου πιστωτικού λογαριασμού και ο συνολικός αριθμός
πιστωτικών λογαριασμών του δανειολήπτη έχουν περιορισμένη αυτόνομη
προβλεπτική ισχύ.
Η μειωμένη σημαντικότητά τους μπορεί εν μέρει να εξηγηθεί και από
επικάλυψη πληροφορίας με σημαντικότερα features.

% -------------------------------------------------------------------
% Σχήμα 4.7 — SHAP Feature Importance (NB03)
% Τοποθέτηση: αμέσως μετά τα αποτελέσματα Global SHAP
% -------------------------------------------------------------------
\begin{figure}[htbp]
    \centering
    \includegraphics[width=0.80\textwidth]{05_shap_feature_importance}
    \caption[SHAP Global Feature Importance — XGBoost]{SHAP Global Feature Importance για το XGBoost. Το mean(|SHAP value|)
             κάθε χαρακτηριστικού εκφράζει τη μέση συνεισφορά του στις
             προβλέψεις. Κορυφαίο χαρακτηριστικό το \texttt{ExternalRiskEstimate},
             ακολουθούμενο από δείκτες πρόσφατης πιστωτικής δραστηριότητας
             και χρέους.}
    \label{fig:shap_importance}
\end{figure}

% 06_shap_summary.png
Το Σχήμα~\ref{fig:shap_summary} παρουσιάζει το Beeswarm διάγραμμα
για το μοντέλο XGBoost.
Το παρόν διάγραμμα προσθέτει κατεύθυνση και κατανομή στο μονοδιάστατο
bar chart του Σχήματος~\ref{fig:shap_importance}.
Κάθε τελεία αντιστοιχεί σε μια εγγραφή του test set και το χρώμα της
δείχνει αν η τιμή του χαρακτηριστικού είναι υψηλή (κόκκινο) ή χαμηλή (μπλε).
Η σειρά των μεταβλητών είναι με βάση τη σημαντικότητά τους ενώ η
συντεταγμένη κάθε τελείας στον άξονα x δείχνει την SHAP τιμή της,
δηλαδή το πόσο και κατά ποιον τρόπο επηρεάζει την τελική πρόβλεψη.
Για το \texttt{ExternalRiskEstimate} παρατηρείται σαφής διαχωρισμός
χρωμάτων στον άξονα x, όπως και για τα \texttt{NetFractionRevolvingBurden},
\texttt{AverageMInFile}, \texttt{PercentTradesNeverDelinq} και
\texttt{NumSatisfactoryTrades}.
Το γεγονός αυτό αποκαλύπτει ότι οι τιμές των αναφερθέντων μεταβλητών
αποτελούν ισχυρό κριτήριο για την τελική συνεισφορά τους στο ποσοστό επικινδυνότητας.
Αντίθετα, για τις μεταβλητές \texttt{MSinceMostRecentInqexcl7days},
\texttt{PercentInstallTrades}, \texttt{NumRevolvingTradesWBalance},
\texttt{MSinceMostRecentDelq} και \texttt{NumTradesOpenInLast12M}
παρατηρείται έντονη ανάμειξη χρωμάτων.
Η τιμή των μεταβλητών δεν καθορίζει την SHAP τιμή τους με αποτέλεσμα
δανειολήπτες με ίδια τιμή στις μεταβλητές να επηρεάζονται διαφορετικά.
Το φαινόμενο έχει τρεις πιθανές εξηγήσεις που για κάθε μεταβλητή
μπορεί να διαφέρει: η αλληλεπίδραση με άλλα χαρακτηριστικά του
δανειολήπτη, η χαμηλή προβλεπτική ισχύς του χαρακτηριστικού που
δεν επιτρέπει την ξεκάθαρη συσχέτισή του με τον κίνδυνο, ή η μη
μονότονη σχέση του με αυτόν.
Τα συμπεράσματα από το διάγραμμα είναι σημαντικά για την
ερμηνευσιμότητα του μοντέλου, καθώς οπτικοποιείται η κατεύθυνση
και η ένταση επίδρασης κάθε χαρακτηριστικού και συμβάλλει στη
διαδικασία επεξηγηματικότητας προς τη συμμόρφωση με το κανονιστικό
πλαίσιο~\citep{eba2023}.

% -------------------------------------------------------------------
% Σχήμα 4.8 — SHAP Beeswarm Plot XGBoost (NB03)
% Τοποθέτηση: αμέσως μετά το κείμενο ανάλυσης
% -------------------------------------------------------------------
\begin{figure}[htbp]
    \centering
    \includegraphics[width=\textwidth]{06_shap_summary}
    \caption[SHAP Beeswarm Plot — XGBoost]{SHAP Beeswarm Plot για το XGBoost. Κάθε τελεία αντιστοιχεί
             σε έναν δανειολήπτη του test set. Το χρώμα δείχνει την τιμή
             του χαρακτηριστικού (κόκκινο = υψηλή, μπλε = χαμηλή) και η
             θέση στον άξονα x την κατεύθυνση και ένταση επίδρασής του
             στην πρόβλεψη.}
    \label{fig:shap_summary}
\end{figure}

\subsection{Ερμηνευσιμότητα SHAP --- Τοπικό Επίπεδο}
\label{sec:results_shap_local}

% 07_shap_local_bad.png, 08_shap_local_good.png
Όπως αναφέρθηκε και σε προηγούμενες παραγράφους ο κανονισμός GDPR \citep{goodman2017european} και οι
κατευθύνσεις της EBA \citep{eba2023} προστατεύουν το «δικαίωμα εξήγησης» των υποψηφίων
δανειοληπτών.
Τα Σχήματα~\ref{fig:shap_local_bad} και~\ref{fig:shap_local_good} ικανοποιούν
συγκεκριμένα αυτό το δικαίωμα.
Παρουσιάζονται δύο SHAP Waterfall plots, ένα για Bad και ένα για Good δανειολήπτη.
Για κάθε δανειολήπτη το διάγραμμα ξεκινά από την τιμή 0{,}0 του άξονα x
ο οποίος συμβολίζει το base value, δηλαδή τη μέση πρόβλεψη του μοντέλου.
Κάθε χαρακτηριστικό μετατοπίζει την πρόβλεψη προς τα δεξιά (αυξάνει τον
κίνδυνο) ή αριστερά (μειώνει τον κίνδυνο).
Η σύγκριση ανάμεσα στα δύο διαγράμματα αποκαλύπτει την συλλογιστική
πορεία του μοντέλου.

Για τον Bad δανειολήπτη φαίνεται ότι η μεταβλητή
\texttt{ExternalRiskEstimate}=58 κυριαρχεί στο αποτέλεσμα μετατοπίζοντας
το αποτέλεσμα κατά απόλυτη τιμή 3 φορές περισσότερο από την δεύτερη
μεγαλύτερη μετατόπιση.
Ενώ οι πέντε από τις εννέα πιο σημαντικές μεταβλητές αφαιρούν από το
σκορ επικινδυνότητας του, η τελική τιμή $f(x) = 1{,}059$ (log-odds)
αντιστοιχεί σε υψηλή πιθανότητα αθέτησης.
Για τον Good δανειολήπτη φαίνεται μια συμφωνία στις ισχυρότερες μεταβλητές
με εφτά από τις εννέα να μειώνουν την τελική τιμή $f(x)$ η οποία τελικά
είναι ίση με $-1{,}678$ (log-odds).
Από το διάγραμμα SHAP feature importance προέκυψε το συμπέρασμα ότι η
μεταβλητή \texttt{ExternalRiskEstimate} έχει την μεγαλύτερη σημαντικότητα
για το μοντέλο αλλά για τον συγκεκριμένο δανειολήπτη φαίνεται να μην
επηρεάζει τόσο το τελικό αποτέλεσμα αλλά ούτε να συμφωνεί κατά πρόσημο
με αυτό.
Το \texttt{ExternalRiskEstimate} για τον Good δανειολήπτη έχει θετική SHAP
τιμή (σπρώχνει προς Bad) παρότι ο δανειολήπτης είναι Good --- αυτό δείχνει
ότι άλλα features αντισταθμίζουν την επίδρασή του.
Επιπλέον, η απόλυτη μετατόπιση που προκαλεί στην $f(x)$ είναι όγδοη κατά
φθίνουσα σειρά ίση με $0{,}11$ σε σύγκριση με το $+1$ για τον Bad.

% -------------------------------------------------------------------
% Σχήμα 4.9 — SHAP Waterfall Bad δανειολήπτης (NB03)
% Τοποθέτηση: αμέσως μετά το κείμενο ανάλυσης
% -------------------------------------------------------------------
\begin{figure}[htbp]
    \centering
    \includegraphics[width=\textwidth]{07_shap_local_bad}
    \caption[SHAP Waterfall — Bad δανειολήπτης]{SHAP Waterfall plot για Bad δανειολήπτη. Το
             \texttt{ExternalRiskEstimate}=58 κυριαρχεί στην πρόβλεψη
             με μετατόπιση τριπλάσια από το δεύτερο feature.
             Τελική πρόβλεψη $f(x)=1{,}059$ (log-odds).}
    \label{fig:shap_local_bad}
\end{figure}

% -------------------------------------------------------------------
% Σχήμα 4.10 — SHAP Waterfall Good δανειολήπτης (NB03)
% Τοποθέτηση: αμέσως μετά το Bad waterfall
% -------------------------------------------------------------------
\begin{figure}[htbp]
    \centering
    \includegraphics[width=\textwidth]{08_shap_local_good}
    \caption[SHAP Waterfall — Good δανειολήπτης]{SHAP Waterfall plot για Good δανειολήπτη. Εφτά από τις
             εννέα κορυφαίες μεταβλητές μειώνουν τον κίνδυνο.
             Τελική πρόβλεψη $f(x)=-1{,}678$ (log-odds).}
    \label{fig:shap_local_good}
\end{figure}


\subsection{Partial Dependence Plots}
\label{sec:results_pdp}

% 09_pdp.png
Το κανονιστικό πλαίσιο του τομέα διακρίνεται σε δύο επίπεδα.
Το πρώτο απαιτεί την απάντηση στο ερώτημα «γιατί έγινε η απόρριψη»
ενώ το δεύτερο στο ερώτημα «πως μπορεί να αλλάξει η απόφαση».
Τα Partial Dependence Plots του Σχήματος~\ref{fig:pdp} απαντούν
στο δεύτερο ερώτημα.
Τα PDP για τα τέσσερα κορυφαία χαρακτηριστικά αποκαλύπτουν πως
αυτά επηρεάζουν τη μέση πρόβλεψη, με τα υπόλοιπα χαρακτηριστικά
να παραμένουν σταθερά στη μέση τιμή τους.
Η παραδοχή αυτή για τα υπόλοιπα χαρακτηριστικά δεν μας παρέχει
την πλήρη εικόνα για την υποστήριξη actionable recourse καθώς
αγνοεί πιθανές αλληλεπιδράσεις ανάμεσα στα χαρακτηριστικά.
Αξίζει να σημειωθεί ότι η μεταβλητή
\texttt{MSinceMostRecentInqexcl7days} είναι η μόνη από τις 4 για
την οποία η τιμή της δεν αποτελεί ισχυρό κριτήριο για την τελική
συνεισφορά της στο ποσοστό επικινδυνότητας
(Σχήμα~\ref{fig:shap_summary}).
Έτσι η διατήρηση των υπόλοιπων μεταβλητών στην μέση τιμή τους
έχει μεγαλύτερη βαρύτητα στην ερμηνεία του διαγράμματος αυτής
της μεταβλητής καθώς όταν οι τιμές τους διαφοροποιηθούν είναι
πιθανόν να το επηρεάσουν.

Στα διαγράμματα βλέπουμε στον άξονα x τις τιμές της κάθε
μεταβλητής ενώ στον άξονα y την εξάρτηση της τελικής πρόβλεψης
από αυτά.
Στο διάγραμμα υπάρχει και rug plot με κατακόρυφα σημάδια στον
άξονα x να οπτικοποιούν την πυκνότητα των πραγματικών δειγμάτων.
Παρατηρείται ότι τα περισσότερα διαγράμματα δεν είναι αυστηρά
μονότονα ούτε γραμμικά αλλά οι πυκνότερες περιοχές αυτών θα
μπορούσαν να προσεγγιστούν με την χρήση γραμμικών, άρα και
αυστηρά μονότονων, συναρτήσεων.
Το γεγονός αυτό εξηγεί τα συγκρίσιμα αποτελέσματα στην απόδοση
της Logistic Regression με το XGBoost
(Πίνακας~\ref{tab:model_metrics}).
Στη διαδικασία προσέγγισης η μεταβλητή με την μεγαλύτερη απόκλιση
θα ήταν η \texttt{MSinceMostRecentInqexcl7days}, εύρημα που
ενισχύει τις υποθέσεις της προηγούμενης παραγράφου.
Συγκεκριμένα, το διάγραμμα της \texttt{ExternalRiskEstimate}
παρουσιάζει κορυφή ίση με $0{,}7$ στον άξονα y, είναι φθίνουσα
στο μεγαλύτερο διάστημά της με τις μικρότερες τιμές της να είναι
κοντά στο $0{,}35$.
Η \texttt{MSinceMostRecentInqexcl7days} ενώ αρχικά είναι αύξουσα,
στη συνέχεια παραμένει σταθερή για τιμές x από $-5$ μέχρι $1$
και έπειτα φθίνει.
Το πεδίο τιμών της είναι $(0{,}35\text{--}0{,}57)$.
Η \texttt{NetFractionRevolvingBurden} είναι κυρίως αύξουσα σε
όλο το πεδίο ορισμού της με ελάχιστη τιμή $0{,}47$ και μέγιστη
$0{,}63$.
Τέλος, η \texttt{AverageMInFile} παρουσιάζει φθίνουσα τάση και
παίρνει τιμές ανάμεσα στο $0{,}62$ και $0{,}46$.
Τα πεδία τιμών των μεταβλητών είναι παρόμοια, γεγονός που
επιβεβαιώνει τη συγκρίσιμη κατά απόλυτη τιμή επίδρασή τους
στην τελική πρόβλεψη.

% -------------------------------------------------------------------
% Σχήμα 4.11 — Partial Dependence Plots top-4 features XGBoost (NB03)
% Τοποθέτηση: αμέσως μετά το κείμενο ανάλυσης
% -------------------------------------------------------------------
\begin{figure}[htbp]
    \centering
    \includegraphics[width=\textwidth]{09_pdp}
    \caption[Partial Dependence Plots — top-4 χαρακτηριστικά]{Partial Dependence Plots για τα τέσσερα κορυφαία
             χαρακτηριστικά του XGBoost. Ο άξονας y δείχνει τη μέση
             προβλεπόμενη πιθανότητα Bad. Τα κατακόρυφα σημάδια στον
             άξονα x (rug plot) οπτικοποιούν την πυκνότητα των
             πραγματικών δειγμάτων.}
    \label{fig:pdp}
\end{figure}


\subsection{Ερμηνευσιμότητα SHAP για FFNN και Υπολογιστικό Κόστος}
\label{sec:results_ffnn_shap}

% 10_shap_ffnn_importance.png
Το Σχήμα~\ref{fig:shap_ffnn_importance} παρουσιάζει το SHAP Feature Importance
για το FFNN, υπολογισμένο με τον KernelExplainer.
Σημαντικότερο χαρακτηριστικό αναδεικνύεται το \texttt{ExternalRiskEstimate}
όπως και στο XGBoost (Σχήμα~\ref{fig:shap_importance}).
Τα δύο μοντέλα έχουν συγκρίσιμη κατάταξη καθώς τα υπόλοιπα χαρακτηριστικά
ακολουθούν παρόμοια σειρά με εξαίρεση το
\texttt{MSinceMostRecentInqexcl7days} που ενώ για το XGBoost έχει τη
δεύτερη μεγαλύτερη σημαντικότητα, για το FFNN βρίσκεται στην έκτη θέση.
Αξίζει να σημειωθεί ότι τα δύο διαγράμματα, ενώ υπολογίζουν τη σημαντικότητα
της κάθε μεταβλητής, χρησιμοποιούν διαφορετική κλίμακα στον άξονα x.
Αυτό συμβαίνει διότι το XGBoost επιστρέφει τιμές σε log-odds space
(raw scores) --- οι τιμές που υπολογίζει το SHAP είναι σε αυτή την κλίμακα.
Αντίθετα, το FFNN επιστρέφει απευθείας πιθανότητες ($0$--$1$), άρα οι
SHAP τιμές είναι περιορισμένες σε ένα αρκετά στενότερο εύρος, οπότε οι
τιμές τους δεν είναι άμεσα συγκρίσιμες.
Επιπλέον, η ανάλυση των τιμών SHAP για το FFNN έχει γίνει μέσω
KernelExplainer, το οποίο κάνει προσέγγιση σε περιορισμένο δείγμα σε
αντίθεση με τον TreeExplainer.
Επιπλέον αυτό οδηγεί σε αραίωση των τιμών, δίνοντάς τους πιο ομοιόμορφη κατανομή.
Σημαντική διαφορά ανάμεσα στις δύο συναρτήσεις εξήγησης είναι ο
απαιτούμενος χρόνος υπολογισμού.
Το υπολογιστικό κόστος του KernelExplainer είναι τάξεις μεγαλύτερο, αφού
απαίτησε ${\sim}2{,}5$ ώρες υπολογισμού έναντι δευτερολέπτων για τον
TreeExplainer --- διαφορά που αναλύεται στη συνέχεια.

% -------------------------------------------------------------------
% Σχήμα 4.12 — SHAP Feature Importance FFNN (NB03)
% Τοποθέτηση: αμέσως μετά το κείμενο ανάλυσης
% -------------------------------------------------------------------
\begin{figure}[htbp]
    \centering
    \includegraphics[width=0.80\textwidth]{10_shap_ffnn_importance}
    \caption[SHAP Feature Importance — FFNN]{SHAP Feature Importance για το FFNN υπολογισμένο με
             KernelExplainer. Κορυφαίο χαρακτηριστικό το
             \texttt{ExternalRiskEstimate}, συνεπές με το αντίστοιχο
             αποτέλεσμα του XGBoost. Οι τιμές του άξονα x είναι μη
             συγκρίσιμες με αυτές του Σχήματος~\ref{fig:shap_importance}
             λόγω διαφορετικής κλίμακας εξόδου των δύο μοντέλων.}
    \label{fig:shap_ffnn_importance}
\end{figure}



% 11_shap_comparison.png
Το Σχήμα~\ref{fig:shap_comparison} συγκρίνει τα διαγράμματα SHAP importance
των δύο μοντέλων για τα 10 πιο σημαντικά χαρακτηριστικά του XGBoost
(Σχήμα~\ref{fig:shap_importance}).
Αξίζει να σημειωθεί ότι 8 από τα 10 αυτά χαρακτηριστικά ανήκουν επίσης
στα 10 σημαντικότερα του FFNN (Σχήμα~\ref{fig:shap_ffnn_importance}).
Εξαίρεση αποτελούν το \texttt{MaxDelq2PublicRecLast12M} και το
\texttt{NumInqLast6M}, τα οποία κατατάσσονται στη 12η και 14η θέση
αντίστοιχα για το XGBoost αλλά στην 7η και 10η για το FFNN οπότε
δεν εμφανίζονται στο αντίστοιχο panel του διαγράμματος.
Συμπερασματικά τα δύο μοντέλα παρουσιάζουν υψηλή συμφωνία ως προς τη σχετική
βαρύτητα των χαρακτηριστικών.

% -------------------------------------------------------------------
% Σχήμα 4.13 — SHAP Comparison XGBoost vs FFNN (NB03)
% Τοποθέτηση: αμέσως μετά το κείμενο ανάλυσης
% -------------------------------------------------------------------
\begin{figure}[htbp]
    \centering
    \includegraphics[width=\textwidth]{11_shap_comparison}
    \caption[Σύγκριση SHAP XGBoost vs FFNN]{Σύγκριση SHAP Feature Importance XGBoost και FFNN για τα
             10 σημαντικότερα χαρακτηριστικά του XGBoost.Οι δύο άξονες 
             x έχουν διαφορετική κλίμακα λόγω της διαφορετικής
            κλίμακας εξόδου των δύο μοντέλων --- η σύγκριση αφορά τη σχετική
            κατάταξη των χαρακτηριστικών, όχι τα απόλυτα μεγέθη.}
    \label{fig:shap_comparison}
\end{figure}




% 24_kernelexplainer_convergence.png
Το Σχήμα~\ref{fig:kernelexplainer_convergence} παρουσιάζει τα αποτελέσματα
του KernelExplainer για δύο προσεγγίσεις, μικρής και μέτριας υπολογιστικής
απαίτησης, και για την πλήρη ανάλυση.
Στην αριστερή πλευρά υπάρχει το διάγραμμα κανονικοποιημένης σημαντικότητας
των 10 σημαντικότερων χαρακτηριστικών του FFNN.
Στη δεξιά πλευρά εμφανίζεται το διάγραμμα κατάταξης συσχέτισης Spearman
για τις δύο προσεγγίσεις.

Στο διάγραμμα κανονικοποιημένης σημαντικότητας παρατηρείται ότι ενώ οι
τιμές των δύο προσεγγίσεων δεν διαφέρουν σημαντικά από αυτές της πλήρους
ανάλυσης, η διαφορά τους θα είχε ως αποτέλεσμα τη διαφορετική κατάταξη
σημαντικότητας των μεταβλητών.
Συγκεκριμένα, για τη μεταβλητή \texttt{NumSatisfactoryTrades} η μικρότερη
υπολογιστικά προσέγγιση έχει αποδώσει μικρότερη σημαντικότητα από την
\texttt{MSinceMostRecentInqexcl7days}, ενώ στην πλήρη ανάλυση το
\texttt{NumSatisfactoryTrades} προηγείται στην κατάταξη.
Παρόμοια παρατήρηση γίνεται για τις μεταβλητές \texttt{PercentInstallTrades}
και \texttt{NumRevolvingTradesWBalance} για τη μεγαλύτερη υπολογιστικά
προσέγγιση.
Οι παρατηρήσεις αυτές υπογραμμίζουν το γεγονός ότι χωρίς την πλήρη
ανάλυση, η οποία είναι δυσανάλογα απαιτητική σε υπολογιστικούς πόρους,
δεν μπορούμε να είμαστε σίγουροι για τα αποτελέσματα.

Στο διάγραμμα κατάταξης συσχέτισης Spearman παρατηρούμε ότι και οι δύο
προσεγγίσεις ξεπερνούν τόσο το όριο $0{,}90$ της αποδεκτής συσχέτισης
όσο και το όριο υψηλής συσχέτισης που είναι ίσο με $0{,}95$.
Οι τιμές τους $0{,}986$ και $0{,}992$ είναι αρκετά υψηλές και αποδεικνύουν
μια ακριβή προσέγγιση του αποτελέσματος της πλήρους ανάλυσης.
Παρά την διαφορά στην κατάταξη των χαρακτηριστικών, συνολικά τα
αποτελέσματα δεν διαφέρουν σημαντικά.
Εμπειρικά αποδεικνύεται ότι για το παρόν μοντέλο και βάση δεδομένων θα
μπορούσε να γίνει μια από τις δύο προσεγγίσεις εξοικονομώντας αρκετό
χρόνο, εφόσον οι προσεγγίσεις είναι της τάξης των μερικών λεπτών ενώ η
πλήρης προσέγγιση χρειάζεται $2{,}5$ ώρες.
Το αποτέλεσμα αυτό δεν μπορεί να γενικευτεί για άλλες βάσεις δεδομένων
και μοντέλα, καθώς χωρίς την πλήρη ανάλυση δεν γίνεται να είμαστε
σίγουροι για τον ρυθμό σύγκλισης των προσεγγίσεων.
Τέλος, αξίζει να αναφερθεί η σημαντική διαφορά της υπολογιστικής
απαίτησης ανάμεσα στα μοντέλα XGBoost και FFNN, δηλαδή ανάμεσα στις
συναρτήσεις TreeExplainer και KernelExplainer.
Η πλήρης ανάλυση του TreeExplainer διαρκεί μόλις λίγα δευτερόλεπτα σε
αντίθεση με τις $2{,}5$ ώρες για τον KernelExplainer.
Η ανάλυση για το μοντέλο XGBoost είναι τάξεις μικρότερη, έτσι ώστε
ακόμα και η διαφορά της με τις προσεγγίσεις των 2 και 10 λεπτών να
είναι εμφανής.
Η διαφορά αυτή προφανώς αυξάνεται και γίνεται δυσκολότερα διαχειρίσιμη
για μεγαλύτερες βάσεις δεδομένων.

% -------------------------------------------------------------------
% Σχήμα 4.14 — KernelExplainer Convergence (NB03)
% Τοποθέτηση: αμέσως μετά το κείμενο ανάλυσης
% -------------------------------------------------------------------
\begin{figure}[htbp]
    \centering
    \includegraphics[width=\textwidth]{24_kernelexplainer_convergence}
    \caption[Σύγκλιση KernelExplainer — FFNN]{Σύγκλιση KernelExplainer για το FFNN. Αριστερά: κανονικοποιημένη
             σημαντικότητα SHAP για τρεις διαμορφώσεις (${\sim}2$ λεπτά,
             ${\sim}10$ λεπτά, ${\sim}2{,}5$ ώρες). Δεξιά: Spearman rank
             correlation κάθε προσέγγισης έναντι της πλήρους ανάλυσης.
             Και οι δύο προσεγγίσεις ξεπερνούν το όριο υψηλής συσχέτισης
             $0{,}95$, αποδεικνύοντας εμπειρικά αποδεκτή σύγκλιση για
             το παρόν dataset.}
    \label{fig:kernelexplainer_convergence}
\end{figure}

%==============================================================
\section{Ανίχνευση και Μετριασμός Μεροληψίας}
\label{sec:results_bias}

\subsection{Ανίχνευση Μεροληψίας}
\label{subsec:bias_detection}

Η ανίχνευση μεροληψίας πραγματοποιήθηκε και στα τρία μοντέλα με τις μετρικές
Δημογραφικής Ισότητας (DPD) και Ανισομερούς Επίπτωσης (DI) που περιεγράφηκαν
στην ενότητα~\ref{sec:bias_methodology}. Οι δανειολήπτες χωρίστηκαν στις ομάδες Α και Β
ανάλογα με την τιμή του χαρακτηριστικού \texttt{AverageMInFile} το οποίο
συμβολίζει την παλαιότητα του πιστωτικού τους ιστορικού. Η Ομάδα Α αντιστοιχεί
σε \texttt{AverageMInFile} κάτω της διαμέσου (76 μήνες) (βραχύ πιστωτικό
ιστορικό) και η Ομάδα Β σε τιμές πάνω της διαμέσου (παλαιό πιστωτικό ιστορικό). Ο διαχωρισμός βάση της διαμέσου 
επιτρέπει τη δημιουργία ομάδων ίσου μεγέθους διασφαλίζοντας έτσι τη στατιστική αξιοπιστία των group fairness metrics \citep{barocas2019fairness}.
Η επιλογή της συγκεκριμένης μεταβλητής εξηγήθηκε αναλυτικά στην
ενότητα~\ref{sec:bias_methodology}.

Ο Πίνακας~\ref{tab:bias_results} περιέχει τα ποσοστά έγκρισης της κάθε ομάδας
καθώς και τους δείκτες DPD και DI που υπολογίζονται από αυτά για τα τρία μοντέλα.
Το ιδανικό αποτέλεσμα θα ήταν οι δύο ομάδες να έχουν παρόμοια ποσοστά έγκρισης
και δείκτες DPD κάτω του 0{,}1 και DI άνω του 0{,}8. Τέτοιο αποτέλεσμα δεν
παρατηρείται σε κανένα μοντέλο. Και οι τρεις αλγόριθμοι έχουν σχεδόν ίδια
ποσοστά έγκρισης, 30\% για την Α και 62\% για την Β, ενώ οι διαφορές στους
δύο δείκτες είναι αμελητέες. Στον δείκτη DPD τα μοντέλα έχουν βαθμολογία
${\sim}0{,}33$ και στον DI ${\sim}0{,}47$, και οι δύο πολύ μακριά των αποδεκτών
ορίων. Παρατηρείται λοιπόν μεροληψία και στους τρεις αλγορίθμους υπέρ των
δανειοληπτών της Ομάδας Β.

% -------------------------------------------------------------------
% Πίνακας 4.3 — Bias Detection αποτελέσματα
% Τοποθέτηση: αμέσως μετά την πρώτη παράγραφο
% -------------------------------------------------------------------
\begin{table}[htbp]
    \centering
    \caption[Αποτελέσματα ανίχνευσης μεροληψίας]{Αποτελέσματα ανίχνευσης μεροληψίας για τα τρία μοντέλα.
    Όριο αποδοχής: DPD $< 0{,}10$, DI $\geq 0{,}80$.}
    \label{tab:bias_results}
    \begin{tabular}{lcccc}
        \toprule
        \textbf{Μοντέλο} & \textbf{Ομάδα Α} & \textbf{Ομάδα Β}
        & \textbf{DPD} & \textbf{DI} \\
        \midrule
        Logistic Regression & 30{,}1\% & 63{,}9\% & 0{,}3381 & 0{,}4712 \\
        XGBoost             & 29{,}6\% & 62{,}4\% & 0{,}3287 & 0{,}4737 \\
        FFNN                & 30{,}4\% & 62{,}7\% & 0{,}3233 & 0{,}4846 \\
        \bottomrule
    \end{tabular}
\end{table}

Το Σχήμα~\ref{fig:bias_comparison} παρουσιάζει οπτικά τα αποτελέσματα του
Πίνακα~\ref{tab:bias_results}. Το γεγονός ότι και τα τρία μοντέλα εμφανίζουν
παρόμοια στατιστικά και βαθμολογία στους δείκτες μας οδηγεί στο συμπέρασμα
ότι η μεροληψία δημιουργείται από την βάση δεδομένων \citep{mehrabi2021survey, suresh2021framework}. 
Η εμφάνιση μεροληψίας
στην συγκεκριμένη προσομοίωση δεν υποδεικνύει λάθος στην «εσωτερική λογική»
κάποιου μοντέλου καθώς όλα τους παράγουν ίδια αποτελέσματα. Η μόνη λογική
εξήγηση είναι ότι η βάση δεδομένων περιέχει δεδομένα πιστωτικού κίνδυνου τα
οποία ευνοούν τους δανειολήπτες με μακρύ πιστωτικό ιστορικό. Το γεγονός αυτό
δεν καθιστά προβληματική την βάση καθώς το συγκεκριμένο χαρακτηριστικό μπορεί
να έχει πραγματικό αντίκτυπο στην πιθανότητα αθέτησης του δανείου. Επιπλέον,
υπάρχει η πιθανότητα η βάση δεδομένων να μην είναι αντιπροσωπευτική και τελικά
αυτός να είναι ο λόγος που παράγει μεροληψία. Και στις δύο περιπτώσεις οι
ρυθμιστικές αρχές απαιτούν έλεγχο και αντιμετώπιση της μεροληψίας \citep{eba2023} ακόμα και
εάν είναι εις βάρος της απόδοσης των μοντέλων.

% -------------------------------------------------------------------
% Σχήμα 4.15 — Bias Comparison τριών μοντέλων (NB04)
% Τοποθέτηση: αμέσως μετά την τελευταία παράγραφο
% -------------------------------------------------------------------
\begin{figure}[htbp]
    \centering
    \includegraphics[width=\textwidth]{12_bias_comparison}
    \caption[Σύγκριση μετρικών μεροληψίας τριών μοντέλων]{Σύγκριση μετρικών μεροληψίας για τα τρία μοντέλα. Αριστερά:
             ποσοστά έγκρισης Ομάδας Α (βραχύ ιστορικό) και Ομάδας Β
             (παλαιό ιστορικό). Κέντρο: Demographic Parity Difference ---
             η διακεκομμένη γραμμή δείχνει το όριο $0{,}10$. Δεξιά:
             Disparate Impact Ratio --- η διακεκομμένη γραμμή δείχνει
             το όριο $0{,}80$. Και τα τρία μοντέλα εμφανίζουν DI
             ${\sim}0{,}47$, σημαντικά κάτω του αποδεκτού ορίου.}
    \label{fig:bias_comparison}
\end{figure}


\subsection{Μετριασμός Μεροληψίας}
\label{subsec:bias_mitigation}

Για την επίδειξη της τεχνικής μετριασμού εφαρμόστηκε ο αλγόριθμος
\texttt{ExponentiatedGradient} της βιβλιοθήκης Fairlearn με base estimator τη
Logistic Regression και ο χειροκίνητος ThresholdOptimizer μέσω grid search
για το XGBoost, για τους λόγους που αναφέρθηκαν στην
ενότητα~\ref{sec:bias_methodology}.

Στο Σχήμα~\ref{fig:fairness_mitigation} παρουσιάζονται τα αποτελέσματα για την
Logistic Regression. Στο αριστερό μέρος το διάγραμμα για την μετρική DI πριν
και μετά την διαδικασία συμμόρφωσης αποδεικνύει ότι άλλαξε η συμπεριφορά του
μοντέλου πετυχαίνοντας απόδοση $0{,}931$ από $0{,}471$, μεγαλύτερη του ορίου
$0{,}8$. Στο δεξί μέρος του σχήματος το διάγραμμα της μετρικής AUC
οπτικοποιεί τις επιπτώσεις του mitigation στην απόδοση. Το κόστος του Fairness
mitigation φαίνεται να είναι μια σχετική μείωση της τάξης του 14\% στην AUC
καθώς η βαθμολογία του αλγορίθμου μειώθηκε από $0{,}789$ σε $0{,}680$.

% -------------------------------------------------------------------
% Σχήμα 4.16 — Fairness Mitigation LR (NB04)
% Τοποθέτηση: αμέσως μετά την παράγραφο LR
% -------------------------------------------------------------------
\begin{figure}[htbp]
    \centering
    \includegraphics[width=\textwidth]{20_fairness_mitigation}
    \caption[Fairness mitigation — Logistic Regression]{Αποτελέσματα fairness mitigation για την Logistic Regression
             με \texttt{ExponentiatedGradient}. Αριστερά: DI πριν
             ($0{,}471$) και μετά ($0{,}931$) --- η διακεκομμένη γραμμή
             δείχνει το όριο $0{,}80$. Δεξιά: AUC πριν ($0{,}789$) και
             μετά ($0{,}680$) --- κόστος $0{,}109$ μονάδων.}
    \label{fig:fairness_mitigation}
\end{figure}

Στο Σχήμα~\ref{fig:fairness_mitigation_comparison} συγκρίνονται τα αποτελέσματα
του μοντέλου LR με το XGBoost. Στο αριστερό μέρος, στο διάγραμμα για την
μετρική DI παρατηρείται ότι ο αλγόριθμος XGBoost καταφέρνει τέλεια βαθμολογία
$1{,}000$. Το ίδιο δεν συμβαίνει για τον αλγόριθμο LR διότι γίνεται αυτόματα
με την συνάρτηση \texttt{ExponentiatedGradient} με $\varepsilon = 0{,}01$ η
οποία προσεγγίζει το όριο χωρίς να το εγγυάται. Αξίζει να σημειωθεί ότι για
το XGBoost δεν έγιναν in-processing αλλαγές, δηλαδή επανεκπαίδευση μοντέλου,
αλλά post-processing αλλαγές στο κατώφλι πιθανότητας. Ως αποτέλεσμα η μετρική
AUC παραμένει αμετάβλητη και ίση με $0{,}796$. Όπως είδαμε στο
Σχήμα~\ref{fig:fairness_mitigation} το ίδιο δεν ισχύει για την LR. Τέλος, στο
δεξί μέρος το διάγραμμα Balanced Accuracy παρουσιάζει το κόστος της διαδικασίας
μετριασμού στα δύο μοντέλα. Παρόλο που η μετρική AUC παρέμεινε σταθερή υπάρχει
υπολογίσιμο κόστος και για το μοντέλο XGBoost. Τα αποτελέσματα δίνουν ξεκάθαρο
πλεονέκτημα στο μοντέλο XGBoost διότι κατάφερε τέλεια βαθμολόγηση στην μετρική
DI ενώ δεν χρειάστηκε επανεκπαίδευση η οποία απαιτεί περισσότερους
υπολογιστικούς πόρους.

% -------------------------------------------------------------------
% Σχήμα 4.17 — Fairness Mitigation Comparison LR vs XGBoost (NB04)
% Τοποθέτηση: αμέσως μετά την παράγραφο σύγκρισης
% -------------------------------------------------------------------
\begin{figure}[htbp]
    \centering
    \includegraphics[width=\textwidth]{25_fairness_mitigation_comparison}
    \caption[Σύγκριση fairness mitigation LR vs XGBoost]{Σύγκριση fairness mitigation LR (in-processing,
             \texttt{ExponentiatedGradient}) και XGBoost (post-processing,
             Manual ThresholdOptimizer). Αριστερά: DI πριν και μετά για
             κάθε μοντέλο --- το XGBoost επιτυγχάνει DI $= 1{,}000$.
             Δεξιά: Balanced Accuracy πριν και μετά --- το AUC του
             XGBoost παραμένει αμετάβλητο στο $0{,}796$.}
    \label{fig:fairness_mitigation_comparison}
\end{figure}

Το παρατηρούμενο fairness-accuracy trade-off, όπως και η σχετικά καλύτερη 
συμπεριφορά του post-processing έναντι του in-processing, ευθυγραμμίζονται 
με τα εμπειρικά ευρήματα της βιβλιογραφίας στο fair credit scoring 
\citep{kozodoi2022fairness}.

%==============================================================
\section{Παρακολούθηση Σταθερότητας}
\label{sec:results_monitoring}
\subsection{Σταθερότητα Πληθυσμού --- PSI}
\label{subsec:results_psi}

% 17_psi_scenarios.png
Ο δείκτης PSI υπολογίστηκε για τα τρία σενάρια οικονομικής επιδείνωσης που
περιγράφηκαν στην ενότητα~\ref{sec:psi}. Τα αποτελέσματα για το XGBoost
παρουσιάζονται στο Σχήμα~\ref{fig:psi_scenarios}. Το σενάριο Βασικής Γραμμής
(Baseline) έχει $\text{PSI} = 0{,}000$. Για ήπια επιδείνωση υπολογίστηκε PSI
ίσο με $0{,}069$, που βρίσκεται κάτω από το όριο $0{,}10$ οπότε θεωρείται
σταθερό. Για μέτρια επιδείνωση το PSI είναι ίσο με $0{,}210$, βρίσκεται
ανάμεσα στα όρια $0{,}10$ και $0{,}25$, που υποδηλώνει ανάγκη παρακολούθησης.
Για σοβαρή επιδείνωση $\text{PSI} = 0{,}641$, πολύ πάνω από το $0{,}25$ οπότε
απαιτείται επανεκπαίδευση μοντέλου.

Η μετάβαση μεταξύ των σεναρίων τονίζει την χρησιμότητα του δείκτη PSI.
Υποθέτοντας ότι η επιδείνωση της οικονομίας δεν είναι ακαριαία και έτσι το
μοντέλο θα περνούσε από όλα τα προηγούμενα στάδια πριν καταλήξει σε κάποια
κατάσταση, συμπεραίνουμε ότι μέσω του δείκτη θα λάβουμε μια προειδοποίηση
για την απόδοση του μοντέλου πριν αυτή υποβαθμιστεί αισθητά και οδηγήσει
σε λάθος αποφάσεις στον πιστωτικό κίνδυνο.

% -------------------------------------------------------------------
% Σχήμα 4.18 — PSI Scenarios XGBoost (NB05)
% Τοποθέτηση: αμέσως μετά το κείμενο ανάλυσης
% -------------------------------------------------------------------
\begin{figure}[htbp]
    \centering
    \includegraphics[width=\textwidth]{17_psi_scenarios}
    \caption[PSI σενάρια οικονομικής επιδείνωσης — XGBoost]{PSI για τέσσερα σενάρια οικονομικής επιδείνωσης (XGBoost).
             Χρωματική κωδικοποίηση: πράσινο $< 0{,}10$ (σταθερό),
             πορτοκαλί $0{,}10$--$0{,}25$ (παρακολούθηση),
             κόκκινο $> 0{,}25$ (επανεκπαίδευση).}
    \label{fig:psi_scenarios}
\end{figure}


\subsection{Σταθερότητα Χαρακτηριστικών — CSI}
\label{subsec:results_csi}
% 13_csi.png
Για κάθε χαρακτηριστικό της βάσης FICO HELOC υπολογίστηκε ο δείκτης CSI
στο σενάριο της μέτριας επιδείνωσης. Τα αποτελέσματα παρουσιάζονται στο
Σχήμα~\ref{fig:csi}. Όπως ήταν αναμενόμενο, τα τρία χαρακτηριστικά που
υπέστησαν τις μεγαλύτερες αλλαγές μέσω της προσομοίωσης εμφανίζονται με
κόκκινο χρώμα --- \texttt{ExternalRiskEstimate} και
\texttt{NetFractionRevolvingBurden} έχουν περάσει το όριο αστάθειας ---
και πορτοκαλί χρώμα το \texttt{AverageMInFile}, εφόσον βρίσκεται στην
περιοχή που δηλώνει ανάγκη παρακολούθησης.
Εκτός από τις τρεις αυτές μεταβλητές, στις υπόλοιπες έχει προστεθεί μόνο
λευκός θόρυβος ($0{,}1$ std) ο οποίος δεν προκαλεί σοβαρή μεταβολή στις
τιμές τους, με αποτέλεσμα ο δείκτης CSI τους να παραμένει κάτω από $0{,}10$.
Εξαίρεση αποτελούν τα χαρακτηριστικά \texttt{MaxDelq2PublicRecLast12M}
($\text{CSI} = 0{,}920$) και \texttt{MaxDelqEver} ($\text{CSI} = 0{,}906$),
τα οποία βρίσκονται πέραν του ορίου αστάθειας με τιμές μάλιστα υψηλότερες
από τα σκόπιμα τροποποιημένα χαρακτηριστικά. Η τεχνητά υψηλή τιμή τους
οφείλεται στη διακριτή φύση των τιμών τους και εξηγείται αναλυτικά στην
ενότητα~\ref{sec:csi}.

% -------------------------------------------------------------------
% Σχήμα 4.19 — CSI ανά χαρακτηριστικό (NB05)
% Τοποθέτηση: αμέσως μετά το κείμενο ανάλυσης
% -------------------------------------------------------------------
\begin{figure}[htbp]
    \centering
    \includegraphics[width=\textwidth]{13_csi}
    \caption[CSI ανά χαρακτηριστικό — μέτρια επιδείνωση]{CSI για τα 23 χαρακτηριστικά της βάσης FICO HELOC στο σενάριο
             μέτριας επιδείνωσης. Χρωματική κωδικοποίηση: πράσινο $< 0{,}10$
             (σταθερό), πορτοκαλί $0{,}10$--$0{,}25$ (παρακολούθηση),
             κόκκινο $> 0{,}25$ (αστάθεια). Τα \texttt{MaxDelq2PublicRecLast12M}
             και \texttt{MaxDelqEver} εμφανίζουν τεχνητά υψηλό CSI λόγω
             διακριτής φύσης --- βλ.\ ενότητα~\ref{sec:csi}.}
    \label{fig:csi}
\end{figure}


\subsection{Σύγκριση Παρακολούθησης XGBoost και FFNN}
\label{subsec:results_monitoring_comparison}

% 21_monitoring_comparison.png
Στο Σχήμα~\ref{fig:monitoring_comparison} παρουσιάζεται η σύγκριση
παρακολούθησης για τα μοντέλα XGBoost και FFNN. Η δεξιά πλευρά του
σχήματος παρουσιάζει τα 10 πιο ασταθή χαρακτηριστικά σε ραβδοδιάγραμμα.
Οι τιμές CSI των χαρακτηριστικών είναι ίδιες και για τα δύο μοντέλα
καθώς εξαρτώνται μόνο από τις αλλαγές στα χαρακτηριστικά και όχι στην
πρόβλεψη του μοντέλου. Σε αντίθεση με τις τιμές CSI, τα αποτελέσματα
του δείκτη PSI εξαρτώνται από τις προβλέψεις των μοντέλων. Στο αριστερό
μέρος του σχήματος οπτικοποιούνται οι διαφορές για τους δύο αλγορίθμους.
Φαίνεται ότι και οι δύο αλγόριθμοι επηρεάζονται με τον ίδιο τρόπο από
τις αλλαγές, με τιμές PSI $0{,}036$ / $0{,}221$ / $0{,}653$ για το FFNN
έναντι $0{,}069$ / $0{,}210$ / $0{,}641$ για το XGBoost. Η μετρική PSI
προειδοποιεί επιτυχώς τόσο σε μέτρια όσο και σε σοβαρή επιδείνωση της
οικονομίας, ενώ δεν δίνει λανθασμένη ειδοποίηση σε περίπτωση ήπιας
αλλαγής. Παρά τη διαφορά στην πολυπλοκότητα των μοντέλων, η
παρακολούθησή τους με τους παραπάνω δείκτες δεν διαφέρει σημαντικά.
Η κρίσιμη διαφορά παραμένει στο κομμάτι της ερμηνείας, όπως
διαπιστώθηκε στην υποενότητα~\ref{sec:results_ffnn_shap}. Για το XGBoost,
εάν παρατηρηθεί σοβαρή αστάθεια μέσω της PSI, η συνάρτηση TreeExplainer
θα μπορέσει γρήγορα να εντοπίσει ποια χαρακτηριστικά την προκάλεσαν.
Σε αντίθεση με το XGBoost, το FFNN απαιτεί KernelExplainer, δηλαδή ώρες
υπολογισμού χωρίς εγγύηση ακρίβειας --- ένα απαγορευτικό γεγονός, διότι
ένα υψηλό CSI δεν θα μπορεί εύκολα να συνδεθεί με τις αλλαγές στο
αποτέλεσμα.

% -------------------------------------------------------------------
% Σχήμα 4.20 — Monitoring Comparison XGBoost vs FFNN (NB05)
% Τοποθέτηση: αμέσως μετά το κείμενο ανάλυσης
% -------------------------------------------------------------------
\begin{figure}[htbp]
    \centering
    \includegraphics[width=\textwidth]{21_monitoring_comparison}
    \caption[Σύγκριση παρακολούθησης XGBoost vs FFNN]{Σύγκριση παρακολούθησης XGBoost και FFNN. Αριστερά: PSI ανά
             σενάριο οικονομικής επιδείνωσης --- οι τιμές των δύο μοντέλων
             είναι παρόμοιες. Δεξιά: CSI για τα 10 πιο ασταθή
             χαρακτηριστικά --- ταυτόσημο και για τα δύο μοντέλα, καθώς
             εξαρτάται από την είσοδο και όχι από την εσωτερική λογική
             του μοντέλου.}
    \label{fig:monitoring_comparison}
\end{figure}

%==============================================================
\section{Συγκεντρωτική Αξιολόγηση}
\label{sec:results_summary}
% 23_summary_table.png
Τα αποτελέσματα του κεφαλαίου οδηγούν σε τέσσερα κύρια συμπεράσματα
για τα τέσσερα βασικά σημεία της εργασίας. Όσον αφορά την απόδοση,
τα μοντέλα Logistic Regression, XGBoost και FFNN έχουν οριακές διαφορές
με το XGBoost να προηγείται ελαφρώς. Αυτό οφείλεται κυρίως στην
προεπεξεργασία της βάσης δεδομένων και τη γραμμικοποίησή της.
Η ερμηνευσιμότητα είναι το κύριο χαρακτηριστικό που διαφοροποιεί τους
αλγορίθμους XGBoost και FFNN. Ο πρώτος είναι πλήρως ερμηνεύσιμος με
μικρό υπολογιστικό κόστος, ενώ ο δεύτερος απαιτεί δυσανάλογους πόρους,
απαγορευτικούς για χρήση σε μεγάλες βάσεις δεδομένων με πολλά
χαρακτηριστικά. Όσον αφορά τη μεροληψία, και τα τρία μοντέλα εμφάνισαν
καθαρά σημάδια της. Το γεγονός αυτό αποδεικνύει ότι η αιτία ήταν τα
δεδομένα της βάσης. Επιπλέον, καθίσταται σαφές ότι ο μετριασμός της
μεροληψίας κοστίζει στον τομέα της απόδοσης. Τέλος, η παρακολούθηση
των μοντέλων είναι εξίσου εύκολη χωρίς να εξαρτάται από την
πολυπλοκότητά τους, γεγονός που δεν μειώνει τη σημαντικότητα της
ερμηνευσιμότητας, καθώς αυτή συνεχίζει να είναι αναγκαία. Το μοντέλο
XGBoost φαίνεται να είναι η καλύτερη επιλογή στην πρόβλεψη του
πιστωτικού κινδύνου, τόσο λόγω της εύκολης και πλήρους
ερμηνευσιμότητάς του όσο και της ικανότητάς του στην αναγνώριση μη
γραμμικών σχέσεων. Η πρόταση πλαισίου επικύρωσης του
Κεφαλαίου~\ref{ch:framework} είναι απόρροια των παραπάνω συμπερασμάτων.

% -------------------------------------------------------------------
% Σχήμα 4.21 — Summary Table (NB02)
% Τοποθέτηση: αμέσως μετά το κείμενο ανάλυσης
% -------------------------------------------------------------------
\begin{figure}[htbp]
    \centering
    \includegraphics[width=\textwidth]{23_summary_table}
    \caption[Συγκεντρωτικός πίνακας αξιολόγησης μοντέλων]{Συγκεντρωτικός πίνακας αξιολόγησης των τριών μοντέλων.
             Χρωματική κωδικοποίηση: κόκκινο = εκτός ορίου
             (DI $< 0{,}80$, PSI $> 0{,}25$),
             πορτοκαλί = παρακολούθηση (PSI $0{,}10$--$0{,}25$).}
    \label{fig:summary_table}
\end{figure}

